\section{XOR interference, counting, and the limits of measurement}\label{sec:boolean}
\subsection{A well-defined interference-like experiment}
For two CMs, define
\begin{equation}
 I_k(A,B)=A\oplus\rot^kB,\qquad k\in\Z/4\Z.
\end{equation}
This is a useful Boolean cancellation experiment. It is not yet a quantum superposition. Coincident features cancel, since $1\oplus1=0$; two contributions in different positions survive separately. Surviving distinct features should not be described as amplitude doubling.

Let $\wt(A)$ count the ones using ordinary integers. Directly from $a\oplus b=a+b-2ab$ for bits,
\begin{equation}\label{eq:xorweight}
 \wt(I_k)=\wt(A)+\wt(B)-2C_{A,B}(k),\qquad
 C_{A,B}(k)=\wt(A\land\rot^kB).
\end{equation}
This identity is exact and is also the familiar relationship between Hamming distance and overlap of binary supports. The four complete $16\times16$ interference tables are retained in the supplementary source and in \code{data/interference_tables.json}, rather than typeset in the submission manuscript. The table for $k$ is simply the XOR table with its second argument permuted by $\rot^k$; it does not by itself establish a new additive law.

For $r=[R]=E_0\oplus E_3$, the clockwise orbit is
\begin{equation}
 [R]\longrightarrow[L]\longrightarrow[\neg R]\longrightarrow[\neg L]\longrightarrow[R].
\end{equation}
Its self-overlap and XOR-weight profiles are
\begin{equation}
 C_{r,r}=(2,1,0,1),\qquad \wt(r\oplus\rot^kr)=(0,2,4,2).
\end{equation}
Thus dividing the latter by four gives the four values of $\sin^2(\phi/2)$ at $\phi=0,\pi/2,\pi,3\pi/2$. This is a four-sample agreement caused by the support geometry, not a derivation of continuous sinusoidal phase dynamics. Other carriers give different profiles: $E_0$ gives $(0,2,2,2)$.

\subsection{What can be derived about Hamming measurement?}
The original LM formalism gives coefficient extraction,
\begin{equation}
 \bra{X_i}M_{X\Theta Y}\ket{Y_j}=\Theta_{ij}
\end{equation}
\cite[Section 4.1.4]{droncheff}. Counting the four positive elementary results therefore gives $\wt(\Theta)$. The count is an integer-valued aggregate of Boolean observations, not their XOR.

\begin{theorem}[Conditional uniqueness of normalized support measure]\label{thm:measure}
Suppose $\mu$ is a nonnegative real-valued function on the sixteen CMs with $\mu(0)=0$, $\mu(\ones)=1$, rotation invariance, and finite additivity for disjoint supports:
\begin{equation}
 A\land B=0\quad\Longrightarrow\quad
 \mu(A\oplus B)=\mu(A)+\mu(B).
\end{equation}
Then $\mu(A)=\wt(A)/4$ for every CM.
\end{theorem}
\begin{proof}
The rotation acts transitively on the four atoms $E_j$, so their measures are equal to some $p$. Their disjoint sum is $\ones$, hence $4p=1$. Each CM is the disjoint union of its occupied atoms.
\end{proof}
The assumption of rotation invariance is an equal-weight assumption about the four elementary positions. Operationally, $\mu(A)$ is the probability that $f_A(X,Y)$ is true if the input pair is sampled uniformly. With a nonuniform input distribution one obtains a weighted count instead. Boolean logic alone does not select a physical sampling distribution.

It is valid, after embedding bit entries in $\R$, that
\begin{equation}
 \wt(A)=\|A\|_F^2=\tr(A^TA).
\end{equation}
The products and trace in this equation are ordinary real or integer arithmetic. Under AND/XOR arithmetic, the analogous contraction returns parity and can vanish for a nonzero vector. This distinction prevents an apparent squared-norm formula from being mistaken for a derivation of the physical Born rule. Ellerman's set-based model provides a particularly close prior comparison: it also combines symmetric difference with an external counting interpretation of overlaps \cite{ellerman}.

The count itself is implementable using Boolean gates. For four entries $a,b,c,d$, its three binary digits are
\begin{align}
 q_0&=a\oplus b\oplus c\oplus d,\\
 q_1&=ab\oplus ac\oplus ad\oplus bc\oplus bd\oplus cd,\\
 q_2&=abcd,\qquad \wt(A)=q_0+2q_1+4q_2.
\end{align}
All sixteen cases were checked. The last equality interprets a multi-bit register as an integer; it does not make the scalar addition characteristic two. Likewise $\wt(A\otimes B)=\wt(A)\wt(B)$ for binary support arrays. This tensor-count identity is not a substitute for defining a quantum tensor product of amplitudes.

\subsection{Integer correlation retains orientation but is not XOR-bilinear}
Writing the bit coefficients as ordinary integers, define
\begin{equation}\label{eq:correlation}
 C_{A,B}(k)=\sum_{j=0}^3 a_j b_{j-k\bmod4},\qquad
 \mathcal C(A,B)=\sum_{k=0}^3 C_{A,B}(k)g^k\in\Z[C_4].
\end{equation}
With $\overline g=g^{-1}$ this is $A\star\overline B$ in the integral group algebra. Thus
\begin{equation}
 \mathcal C(B,A)=\overline{\mathcal C(A,B)},\quad
 \mathcal C(g^rA,g^sB)=g^{r-s}\mathcal C(A,B).
\end{equation}
The correlation convention is linear in the first argument after integral lifting; conjugating the evaluated form gives the usual Dirac convention. These are genuine conjugation and covariance identities. However, the embedding of Boolean arrays as $0/1$ integer arrays is not additive for XOR: $\mathcal C(A\oplus A,B)=0$, whereas $2\mathcal C(A,B)$ need not vanish. Hence this is not an $\F_2$-sesquilinear inner product. Usual polarization arguments relating norm preservation to inner-product preservation cannot be assumed in that mixed arithmetic.

\section{The finite operator audit and no-go results}\label{sec:finite}
\subsection{The exact search universe}
Let $\alg_2=\F_2[C_4]$ with the convolution product. A two-branch state is $\Psi=(A,B)^T\in\alg_2^2$. We enumerate all $16^4=65{,}536$ matrices
\begin{equation}
 U=\cm{u_{11}}{u_{12}}{u_{21}}{u_{22}},\qquad
 U\Psi=\binom{u_{11}\star A\oplus u_{12}\star B}{u_{21}\star A\oplus u_{22}\star B}.
\end{equation}
Every operator is tested on all $256$ states. These are all two-by-two \emph{$\alg_2$-linear} operators, not all possible binary circuits, nonlinear maps, measurement models, or operators in arbitrary dimensions. As binary matrices they form a structured subset of the $8\times8$ matrices.

Define $W(A,B)=\wt(A)+\wt(B)$ and the weight-four shell
\begin{equation}
 \Sigma=\{(A,B):W(A,B)=4\},\qquad |\Sigma|=\binom84=70.
\end{equation}
Also define the full self-correlation $\mathcal C(A,A)+\mathcal C(B,B)$ and its signed quotient, whose real value is $C_0-C_2$. Full pair-correlation preservation means equality on every pair of states, not merely on self-pairs.

\begin{table}[htbp]
\centering\small
\begin{tabularx}{\textwidth}{@{}Y r@{}}
\toprule
Recomputed property & Count\\
\midrule
All two-by-two $\alg_2$-linear operators & $65{,}536$\\
Invertible operators & $24{,}576$\\
Involutions $U^2=I$ & $640$\\
Involutions sending $(0,r)$ to $(r,r)$ & $4$\\
Preserve $W$ on all $256$ states & $32$\\
Preserve full integer self-correlation on all states & $32$\\
Preserve full integer pair-correlation on all state pairs & $32$\\
Map the weight-four shell into itself & $512$\\
Shell preservers invertible on the full space & $512$\\
Shell-preserving involutions & $96$\\
Distinct actions on the shell / kernel size & $128\,/\,4$\\
Preserve full self-correlation on the shell & $512$\\
Preserve signed self-form globally / on the shell & $2{,}048\,/\,8{,}192$\\
Preserve full signed pair-correlation globally & $32$\\
\bottomrule
\end{tabularx}
\caption{Recomputed classifications. The two shell-preserver sets (weight and full self-correlation) coincide exactly. Data: \code{finite_summary.json}; membership masks: \code{finite_classification.npz}.}
\label{tab:finite}
\end{table}

\subsection{Why the global answer is exactly 32}
\begin{theorem}[Global coefficient-Hamming isometries]\label{thm:hamming}
On $\alg_2^d$, every $\alg_2$-linear transformation preserving the total Hamming weight of the $4d$ binary coefficients on all vectors is a permutation of the $d$ branches followed by an independent cyclic rotation within each branch. The group is $C_4\wr S_d$, of order $4^d d!$. For $d=2$ it consists of the $32$ matrices
\begin{equation}
 \diag(g^a,g^b),\qquad \cm0{g^a}{g^b}0,\qquad a,b\in\{0,1,2,3\}.
\end{equation}
\end{theorem}
\begin{proof}
Regard the map as an $\F_2$-linear map on $4d$ binary coordinates. Each unit vector has weight one, so its image must be a unit vector. Images of two different unit vectors cannot coincide, since then their weight-two sum would map to zero. Thus the map is a coordinate permutation. $\alg_2$-linearity means that it commutes with simultaneous multiplication by $g$, which is a permutation with $d$ disjoint cycles of length four. A commuting permutation sends a cycle to a cycle and chooses a cyclic offset within it. There are $d!$ cycle permutations and $4^d$ offsets. Conversely these maps plainly preserve weight.
\end{proof}
The proof is intentionally elementary, but the general phenomenon belongs next to the classical MacWilliams theory of Hamming isometries and its finite-ring extensions.  MacWilliams-type extension theorems state, under their respective hypotheses, that Hamming isometries extend to monomial transformations; Wood established the finite-Frobenius-ring extension property and Greferath--Schmidt gave a combinatorial treatment \cite{wood1999,greferath2000}.  The weight used here is more specific: it counts the four binary coefficients inside each $\F_2[C_4]$ symbol, rather than merely whether a ring symbol is nonzero.  Thus \cref{thm:hamming} should be read as a specialized coefficient-basis classification obtained by combining ordinary binary Hamming isometry with the commuting $C_4$ action, not as a new general MacWilliams theorem.

Full integer-correlation preservation includes the zero-lag self-correlation, which is Hamming weight. Thus it cannot \emph{enlarge} this global group. Every listed phase-permutation does preserve the full correlation, so the classification remains exactly $32$. Earlier motivation suggesting that preservation of a richer form might enlarge the isometry group was incorrect: adding preservation constraints can only retain or remove operators.

The larger signed-self-form class does not contradict this theorem. That form forgets some information, has null directions on raw arrays, and is not a quadratic form over the original XOR addition. Requiring the complete pair form reduces it to the same $32$ maps in this finite test.

\subsection{The shell is richer, but still does not yield the target fringe}
The shell admits the involutory mixer
\begin{equation}
 K=\cm1{1\oplus g^2}{1\oplus g^2}1,\qquad K^2=I.
\end{equation}
It is not a global Hamming isometry: $(E_0,0)$ has weight one, whereas its image has weight three. All $512$ shell preservers are invertible, and their restrictions form a group of $128$ permutations of $\Sigma$.

The four-element kernel has a transparent explanation. The binary span of the weight-four vectors in eight coordinates is the even-parity subspace. A linear map equal to the identity on that subspace has the form $I+v\epsilon$, where $\epsilon$ is total parity. Commutation with the two four-cycles requires $v$ to be constant on each cycle, giving four choices. In ring notation they are
\begin{equation}
 I+\begin{pmatrix}\alpha\ones&\alpha\ones\\\beta\ones&\beta\ones\end{pmatrix},
 \qquad \alpha,\beta\in\F_2.
\end{equation}
Each is invertible because $\epsilon(v)=0$. This explains the observed ratio $512/4=128$ without assigning a new physical meaning to the shell.

For every $U$ in the shell-preserving set, every $\Psi\in\Sigma$, and every $S_k=\diag(g^k,1)$, we computed $W$ branch by branch after $US_k\Psi$. Only seven sequences of first-branch weights occurred:
\begin{equation}
 (0,0,0,0),\ (1,1,1,1),\ (2,2,2,2),\ (3,3,3,3),\ (4,4,4,4),
\end{equation}
\begin{equation}
 (1,3,1,3),\qquad (3,1,3,1).
\end{equation}
There is no sequence $(0,2,4,2)$ or $(4,2,0,2)$, including cyclic changes of phase origin. The search covers $512\cdot70$ operator/state cases, each at four phases. This is a complete negative result for that experiment, not a prohibition on all finite-field quantum-like theories. Modal quantum theory deliberately uses different measurement assumptions \cite{modal}.

\subsection{Two explicit failures that must not be hidden}
The first attempted conditional rotation used the masks $[R]=\cm1010$ and $[\neg R]=\cm0101$:
\begin{equation}
 F(A)=(A\land[\neg R])\oplus(\rot A\land[R]).
\end{equation}
For $A=\cm abcd$ this is
\begin{equation}\label{eq:badmask}
 F(A)=\cm cbdd.
\end{equation}
It deletes $a$, sends $E_0$ to zero, and cannot satisfy $F^4=I$. It does not rotate one column as a reversible phase gate. The earlier assertion that its order follows from $\rot^4=I$ was false.

A later toy splitter in the convolution algebra was
\begin{equation}
 B=\cm1{g^2}11,\qquad B^2=(1\oplus g^2)I,
 \qquad B^4=0.
\end{equation}
It can generate the desired four output weights from a specially chosen carrier, but is nilpotent rather than unitary. Also, $r\ket0$ has $W/4=1/2$ before splitting, not one. Apparent conservation at the final detector does not establish conservation through the circuit. No scalar normalization repairs $B$: if $s^2(1+g^2)=1$, squaring gives $0=1$.

There is an independent obstruction to a familiar Hadamard identity. Over this commutative ring, similarity preserves trace. But $Z_B=\diag(g^2,1)$ has trace $g^2+1\ne0$, while $X=\cm0110$ has trace zero. Hence no invertible $H$ can satisfy $HZ_BH^{-1}=X$. These failures motivate changing the amplitude algebra explicitly, rather than relabeling the failed maps as quantum gates.
